\section{ماژول slim\_compressor (SLiM Structured Low-Rank Integrated Matrix Compressor)}
\label{sec:module_slim_compressor}

\subsection{معرفی و انگیزه}
این مکانیزم بر مبنای مبانی نظری پیشرفته در ادبیات توجه و فشرده‌سازی \cite{ref_slimcompressor_2024} توسعه یافته است.

کاهش بار محاسباتی و اندازه وزن‌ها در شبکه‌های عصبی عمیق عموماً از طریق دو شاخه مجزای کوانتیزاسیون مقادیر و تنک‌سازی ساختاریافته دنبال می‌شود؛ اما اعمال مجزای هر یک از این روش‌ها بدون در نظر گرفتن برهم‌کنش خطاهای تقریب، افت قابل‌توجهی در توانمندی بازنمایی مدل پدید می‌آورد. 

ماژول \lr{slim\_compressor} (برگرفته از رویکرد \lr{SLiM}) روشی یکپارچه و تک‌مرحله‌ای (\lr{One-Shot}) برای ادغام کوانتیزاسیون با دقت پایین ($b$-بیت)، تنک‌سازی نیمه‌ساختاریافته با الگوی سخت‌افزار-آگاه 2:4 (\lr{2:4 Structured Sparsity})، و جبران خطای باقی‌مانده با استفاده از تقریب کم‌مرتبه مبتنی بر تجزیه مقادیر منفرد (\lr{SVD}) ارائه می‌دهد. این الگوریتم با افزودن یک مؤلفه باقیمانده کم‌مرتبه، اثرات تخریبی ناشی از هرس وزن‌ها و کوانتیزاسیون یکنواخت را مهار می‌سازد.

\subsection{مبانی نظری و فرمول‌بندی ریاضی}
فرآیند فشرده‌سازی لایه‌ها در الگوریتم \lr{SLiM} طی مراحل چهارگانه زیر فرمول‌بندی می‌شود:
\begin{enumerate}
    \item \textbf{کوانتیزاسیون $b$-بیت وزن اولیه:}
    فرض کنید $\mathbf{W}_0 \in \mathbb{R}^{d_{\text{out}} \times d_{\text{in}}}$ ماتریس وزن اصلی باشد. عملگر کوانتیزاسیون یکنواخت با گام مقیاس‌بندی $s$ به فرم زیر اعمال می‌شود:
    \begin{equation}
    s = \frac{\max(|\mathbf{W}_0|)}{2^{b-1} - 1}, \quad \mathbf{W}_q = s \cdot \operatorname{clamp}\left( \left\lfloor \frac{\mathbf{W}_0}{s} \right\rceil, -2^{b-1}, 2^{b-1} - 1 \right)
    \end{equation}
    \item \textbf{اعمال تنک‌سازی نیمه‌ساختاریافته 2:4:}
    روی ماتریس کوانتیده‌شده $\mathbf{W}_q$ تصویرگر الگوی تنک 2:4 اعمال می‌گردد که در هر بلوک چهارتایی از عناصر سطر، دو عنصر با کمترین اندازه مطلق را صفر می‌کند:
    \begin{equation}
    \mathbf{W}_{sp} = \Pi_{2:4}(\mathbf{W}_q, |\mathbf{W}_q|)
    \end{equation}
    \item \textbf{تشکیل ماتریس خطای باقیمانده (\lr{Residual Error}):}
    اختلاف بین ماتریس متراکم اولیه و وزن تنکِ کوانتیده‌شده محاسبه می‌شود:
    \begin{equation}
    \mathbf{R} = \mathbf{W}_0 - \mathbf{W}_{sp}
    \end{equation}
    \item \textbf{تقریب کم‌مرتبه باقیمانده با SVD و بازسازی وزن نهایی:}
    ماتریس خطای باقیمانده $\mathbf{R}$ تحت تجزیه مقادیر منفرد با مرتبه $r$ قرار می‌گیرد:
    \begin{equation}
    \mathbf{R} \approx \mathbf{U}_r \boldsymbol{\Sigma}_r \mathbf{V}_r^\top = \mathbf{L}_r
    \end{equation}
    در نهایت، ماتریس وزن بازسازی‌شده فشرده به صورت مجموع مؤلفه تنک و تقریب کم‌مرتبه حاصل می‌شود:
    \begin{equation}
    \mathbf{W}_{slim} = \mathbf{W}_{sp} + \mathbf{L}_r
    \end{equation}
\end{enumerate}

\paragraph{تحلیل طیفی و نامساوی‌های پایداری خطا}
برای بررسی میزان دقت تقریب ماتریس وزن یا نگاشت توجه در ماژول \lr{slim\_compressor}، تفاضل ماتریس فاکتورگیری‌شده از نگاشت متراکم اصلی بر اساس نامساوی هافدینگ-فون‌نیومان ارزیابی می‌گردد. فرض کنید $\mathbf{W} \in \mathbb{R}^{d_{\text{in}} \times d_{\text{out}}}$ ماتریس وزن هدف باشد که توسط تقریب کم‌مرتبه یا پراکنده با رتبه موثر $r$ جایگزین شده است. کران بالای خطای بازسازی بر مبنای نرم فروبنیوس به صورت زیر تصریح می‌شود:
\begin{equation}
\label{eq:spectral_bound_slim_compressor}
\| \mathbf{W} - \mathbf{W}_{r} \|_F = \sqrt{\sum_{i=r+1}^{\min(d_{\text{in}}, d_{\text{out}})} \sigma_i^2(\mathbf{W})} \le \epsilon \| \mathbf{W} \|_F
\end{equation}
که در آن $\sigma_i(\mathbf{W})$ بیانگر $i$-امین مقدار منفرد مرتب‌شده ماتریس به صورت نزولی است. نسبت حفظ انرژی طیفی $\eta_r = \frac{\sum_{i=1}^r \sigma_i^2}{\sum_{i=1}^{d} \sigma_i^2}$ تضمین می‌کند که با انتخاب آستانه $\tau \ge 0.90$، بیش از نود درصد از پراکندگی واریانس داده‌ها حفظ شده و افت سیگنال در انتشار گرادیان رو به عقب به حداقل ممکن کاهش می‌یابد.

\paragraph{تحلیل مجانبی و پیچیدگی اشتراک حافظه سخت‌افزاری}
از منظر پیچیدگی محاسباتی و انتقال داده‌ها در سلسومراتب حافظه (حافظه سراسری \lr{HBM} به حافظه محلی \lr{SRAM})، عملیات ضرب ماتریسی در این مکانیزم نیازمند تخصیص بافر موقت با ابعاد فاکتورگیری‌شده است. برای پردازش دنباله‌ای به طول $N$ با بعد مخفی $D$، بار محاسباتی هسته از فرم درجه دو به رابطه زیر تقلیل می‌یابد:
\begin{equation}
\label{eq:complexity_slim_compressor}
\mathcal{C}_{\text{slim\_compressor}} = \mathcal{O}\left( B \cdot N \cdot D \cdot r + B \cdot H \cdot N \cdot d_k \right) \ll \mathcal{O}\left( B \cdot H \cdot N^2 \cdot d_k \right)
\end{equation}
این بهبود مرتبه پیچیدگی، با کاهش چشمگیر حجم تبادل تانسورها در گذرگاه‌های ارتباطی کارت گرافیک، مانع از اشباع پهنای باند حافظه (پدیده \lr{Memory-Bound}) در فرآیند استنتاج طولانی‌مدت می‌گردد.

\subsection{معماری و پیاده‌سازی}
پیاده‌سازی این ماژول در فایل \lr{Codes/framework/compression\_engine/slim\_compressor.py} با کلاس اصلی \lr{SLiMCompressor} انجام یافته است. مشخصه‌های مهندسی این پیاده‌سازی عبارتند از:
\begin{itemize}
    \item \textbf{سازوکار بازگشت امنیتی (\lr{Fallback Mechanism}):} برای هر لایه، خطای نسبی بازسازی با معیار نرم فروبنیوس $\|\mathbf{W}_0 - \mathbf{W}_{slim}\|_F / \|\mathbf{W}_0\|_F$ اندازه‌گیری می‌شود. چنانچه این خطا از آستانه \lr{max\_error} (پیش‌فرض 0.5) فراتر رود، لایه به وزن اصلی بازگردانده می‌شود.
    \item \textbf{پشتیبانی از لایه‌های خطی و کانوولوشن یک‌بعدی (\lr{Conv1D}):} تشخیص مطمئن لایه‌های خطی بر مبنای نوع لایه و مشخصه \lr{nf} (مورد استفاده در مدل‌های خانواده GPT) بدون تکیه بر حدسیات ابعادی انجام می‌گیرد.
    \item \textbf{محافظت از لایه‌های بحرانی:} لایه‌های حساس مانند طبقه‌بند (\lr{classifier})، سر مدل زبانی (\lr{lm\_head}) و لایه‌های نرمال‌سازی با فراخوانی تابع سراسری \lr{should\_compress\_linear} محافظت می‌شوند.
\end{itemize}

\subsection{نتایج تجربی}
نتایج حاصل از ارزیابی‌های مستقل روی بنچ‌مارک‌های استاندارد برای ماژول \lr{slim\_compressor} در جدول \ref{tab:benchmark_slim_compressor} خلاصه شده است:

\begin{table}[H]
\centering
\small
\caption{نتایج بنچ‌مارک تجربی فشرده‌ساز چندگانه \lr{slim\_compressor} روی مدل‌های پایه}
\label{tab:benchmark_slim_compressor}
\begin{tabular}{lccccc}
\toprule
\textbf{مدل پایه \lr{(Base Model)}} & \textbf{دقت \lr{(Accuracy)}} & \textbf{شاخص \lr{F1}} & \textbf{حافظه \lr{(Peak VRAM)}} & \textbf{توان \lr{(Throughput)}} & \textbf{وضعیت} \\
\midrule
\lr{DistilBERT (SST-2)} & \lr{90.60\%} & \lr{0.9093} & \lr{367.66 MB} & \lr{216,886.93 tok/s} & \lr{PASS} \\
\lr{DistilBERT (MNLI)} & \lr{81.10\%} & \lr{0.8104} & \lr{456.73 MB} & \lr{119,037.02 tok/s} & \lr{PASS} \\
\lr{RoBERTa (SST-2)} & \lr{91.51\%} & \lr{0.9181} & \lr{513.66 MB} & \lr{14,275.54 tok/s} & \lr{PASS} \\
\lr{GPT-2} & -- & -- & -- & -- & \lr{NOT\_RUN} \\
\lr{Qwen} & -- & -- & -- & -- & \lr{NOT\_RUN} \\
\bottomrule
\end{tabular}
\end{table}

بر پایه داده‌های تجربی، فشرده‌ساز \lr{slim\_compressor} در مدل‌های انکودر متقارن بازدهی و حفظ دقت قابل توجه‌ای نشان داده است؛ به طوری که در وظیفه SST-2 روی مدل \lr{RoBERTa} به دقت \lr{91.51\%} و در مدل \lr{DistilBERT} به دقت \lr{90.60\%} رسیده است. همچنین در وظیفه پیچیده‌تر استنتاج زبانی MNLI روی \lr{DistilBERT} دقت پایدار \lr{81.10\%} ثبت شده است. ارزیابی مدل‌های خودبازگشتی \lr{GPT-2} و \lr{Qwen} در پایگاه داده اعتبارسنجی با وضعیت \lr{NOT\_RUN} به دلیل اولویت‌بندی مدل‌های طبقه‌بندی متن علامت‌گذاری شده است.

\subsection{محدودیت‌ها و تحلیل شکست}
نقاط ضعف و ریسک‌های معماری این فشرده‌ساز عبارتند از:
\begin{itemize}
    \item \textbf{تداخل خطای کوانتیزاسیون و هرس در لایه‌های عمیق:} در صورتی که طیف مقادیر منفرد ماتریس خطا به کندی افت کند، مرتبه کوچک $r$ در تقریب SVD قادر به جبران اثر کوانتیزاسیون نبوده و خطای نسبی از حد مجاز فراتر می‌رود که منجر به فعال‌سازی حالت بازگشت (\lr{fallback}) و عدم فشرده‌سازی واقعی لایه خواهد شد.
    \item \textbf{سربار حافظه‌ای نگهداری مقادیر SVD:} نگهداری همزمان ماتریس‌های تنک و فاکتورهای کم‌مرتبه $\mathbf{U}_r$ و $\mathbf{V}_r$ در صورتی که مرتبه انتخابی بزرگ باشد، مزیت ذخیره‌سازی را خنثی می‌سازد.
\end{itemize}
